\documentclass[12pt,a4paper]{article}

%% ── packages ─────────────────────────────────────────────────────────────────
%% XeLaTeX: fontspec handles Unicode natively.
%% IMPORTANT: polyglossia (which loads bidi) must be loaded LAST, after hyperref
%% and all other packages that bidi is incompatible with when loaded after it.
\usepackage{fontspec}
\setmainfont{Linux Libertine O}
\setmonofont{DejaVu Sans Mono}[Scale=0.88]
\usepackage{geometry}
\geometry{margin=2.5cm}
\usepackage{amsmath,amssymb}
\usepackage{booktabs}
\usepackage{tabularx}
\usepackage{array}
\usepackage{graphicx}
\usepackage{xcolor}
\usepackage{hyperref}
\hypersetup{colorlinks=true, linkcolor=blue!60!black, citecolor=teal, urlcolor=blue!60!black}
\usepackage{setspace}
\onehalfspacing
\usepackage{caption}
\usepackage{subcaption}
\usepackage{listings}
\lstset{
  basicstyle=\ttfamily\small,
  backgroundcolor=\color{gray!8},
  frame=single,
  breaklines=true,
  columns=fullflexible,
}
\usepackage{enumitem}
\usepackage{float}
\usepackage{multirow}
\usepackage{tcolorbox}
\tcbuselibrary{skins}
\usepackage{natbib}
\bibliographystyle{plainnat}
%% polyglossia + bidi go last
\usepackage{polyglossia}
\setmainlanguage{english}
\setotherlanguage{hebrew}
\newfontfamily\hebrewfont{Noto Serif Hebrew}[Script=Hebrew, Scale=1.05]

%% ── convenience commands ─────────────────────────────────────────────────────
\newcommand{\vlm}{\textsc{VLM}}
\newcommand{\simEN}{\ensuremath{\mathrm{sim}_{\text{en}}}}
\newcommand{\simHH}{\ensuremath{\mathrm{sim}_{\text{he--he}}}}
\newcommand{\simLaBSE}{\ensuremath{\mathrm{sim}_{\text{LaBSE}}}}
\newcommand{\csim}{\ensuremath{\mathrm{csim}_{\text{en}}}}
\newcommand{\SACs}{\ensuremath{\Delta_{\text{concept}}^{\text{scalar}}}}
\newcommand{\SACv}{\ensuremath{\Delta_{\text{concept}}^{\text{vec}}}}
\newcommand{\SACvb}{\ensuremath{\Delta_{\text{concept}}^{\text{vec-both}}}}
\newcommand{\SAC}{\SACs}
\newcommand{\bsF}{\ensuremath{F_1^{\text{BERT}}}}
\newcommand{\Qwen}{\textit{Qwen2-VL-7B-Instruct}}

%% ─────────────────────────────────────────────────────────────────────────────
\title{%
  \textbf{Evaluating Systems Thinking in Visual Explanation Tasks}\\[0.4em]
  \large Using a Vision-Language Model as an Expert Reference for Scoring
  Open-Ended Hebrew Responses\\[0.6em]
  \normalsize Project Report --- Preliminary Results
}
\author{Research Group, Systems Thinking Lab}
\date{March 2026 \quad \small\textit{(327 samples, image groups 1--2)}}

%% ─────────────────────────────────────────────────────────────────────────────
\begin{document}
\maketitle
\tableofcontents
\bigskip
\hrule
\bigskip

% ═══════════════════════════════════════════════════════════════════════════════
\section{Introduction and Motivation}
% ═══════════════════════════════════════════════════════════════════════════════

Systems thinking (ST) is a cognitive capacity that enables individuals to perceive
interconnected structures, feedback loops, and emergent dynamics rather than isolated
cause-effect chains~\citep{senge1990fifth}.  Despite growing interest in ST education,
reliable \emph{automated} assessment of ST competencies in open-ended tasks remains
largely unsolved.

Traditional closed-form tests (multiple-choice, Likert scales) are easy to score but
capture only a narrow, surface-level signal.  Richer evidence comes from \emph{free-text
justifications}: when participants are asked \emph{why} they chose a specific image as
reflecting a given systems concept, their verbal reasoning reveals the depth of mental
model activation.  The central challenge is that scoring such justifications at scale
requires either costly human annotation or a reliable automated proxy.

This project proposes and evaluates a pipeline in which a state-of-the-art
\emph{Vision-Language Model} (\vlm{}) -- \Qwen{} -- acts as an expert reference:
it is shown the same images and concept-grounding questions, produces a systems-thinking
interpretation in English, and the cosine similarity between that reference answer and
the participant's translated answer is used as a proxy score.

\subsection{Research Questions}

\begin{enumerate}[label=\textbf{RQ\arabic*.}]
  \item Does the \vlm{}-based similarity score capture meaningful variation in response
        quality, or does it merely reflect surface lexical similarity to the concept label?
  \item Does response length (word count) correlate with higher similarity scores, and
        is this a genuine proxy of deeper engagement or simply verbosity?
  \item Are there systematic differences across the ten ST concepts in how well
        participants' answers align with expert \vlm{} interpretations?
  \item How do different embedding spaces (multilingual He--He, English En--En,
        cross-lingual LaBSE) compare as scoring bases?
\end{enumerate}

% ═══════════════════════════════════════════════════════════════════════════════
\section{Dataset}
% ═══════════════════════════════════════════════════════════════════════════════

\subsection{Study Design}

Participants were presented with a questionnaire containing ten systems-thinking
concepts expressed as Hebrew sentences (Table~\ref{tab:concepts}).  For each concept,
they selected one of four candidate images that they perceived as best representing the
concept, and then wrote a free-text Hebrew justification explaining their choice.

\begin{table}[H]
\centering
\caption{Ten systems-thinking concept sentences used in the study.}
\label{tab:concepts}
\small
\begin{tabular}{clp{4cm}p{4.8cm}}
\toprule
\textbf{Q\#} & \textbf{English concept} & \textbf{Phonetic transliteration} & \textbf{Hebrew concept phrase} \\
\midrule
Q1  & Work completed successfully
    & \textit{ha-avoda hushlema be-hatzlacha}
    & \texthebrew{העבודה הושלמה בהצלחה} \\[3pt]
Q2  & The decision has been made
    & \textit{ha-hachlatah hitkabla}
    & \texthebrew{ההחלטה התקבלה} \\[3pt]
Q3  & The child has learned an important lesson
    & \textit{ha-yeled lamad shi'ur hashuv}
    & \texthebrew{הילד למד שיעור חשוב} \\[3pt]
Q4  & The ecosystem is functioning well
    & \textit{ha-ma'arechet ha-ekologit mitpakedet hetev}
    & \texthebrew{המערכת האקולוגית מתפקדת היטב} \\[3pt]
Q5  & Similar systems
    & \textit{ma'arachot domot}
    & \texthebrew{מערכות דומות} \\[3pt]
Q6  & The group advanced towards the goal
    & \textit{ha-kvutza hitzlicha lehitkadem le'ever ha-ya'ad}
    & \texthebrew{הקבוצה הצליחה להתקדם לעבר היעד} \\[3pt]
Q7  & The way I understand reality affects what I see
    & \textit{ha-derech she-bah ani mevin et ha-metziut mashpi'ah al ma she-ani ro'eh}
    & \texthebrew{הדרך שבה אני מבין את המציאות משפיעה על מה שאני רואה} \\[3pt]
Q8  & When I try to understand what is happening I sometimes see things differently
    & \textit{kshe-ani menase lehavin et ma she-kore lif'amim ani ro'eh et ha-dvarim acheret}
    & \texthebrew{כשאני מנסה להבין את מה שקורה לפעמים אני רואה את הדברים אחרת} \\[3pt]
Q9  & Man affects the world around him
    & \textit{ha-adam mashpia al ha-olam saviv}
    & \texthebrew{האדם משפיע על העולם סביבו} \\[3pt]
Q10 & Change is made
    & \textit{ha-shinui na'ase}
    & \texthebrew{השינוי נעשה} \\
\bottomrule
\end{tabular}
\end{table}

\subsection{Data Collection and Structure}

Raw questionnaire data was parsed into a structured dataset organised as follows:

\begin{itemize}[noitemsep]
  \item \textbf{questions/} -- 10 Hebrew question files (\texttt{1.txt}--\texttt{10.txt}).
  \item \textbf{images/} -- 20 images: each question has 2 associated images\\
        (\texttt{<q>\_{1}.png} and \texttt{<q>\_{2}.png}).
  \item \textbf{answers/} -- One file per (question, image-group, participant) triple,
        named \texttt{<question>\_<image>\_<participant>.txt}.
\end{itemize}

The current dataset contains \textbf{539 answer files} across \textbf{10 questions},
\textbf{4 image groups} (answer\_num 1--4), and variable participant counts per
question/image combination.  A missing-pairs audit identified 167 (subject, question,
answer) triples where participants selected an image option but provided no free-text
justification (\texttt{missing\_text}), plus a smaller number of entries where no option
was selected at all (\texttt{missing\_option}).

\medskip
\begin{tcolorbox}[title=Data Completeness Note,colback=yellow!8,colframe=orange!60!black]
Of the 539 answer files present, \textbf{327 are fully scoreable} (image groups 1--2,
which have corresponding image files in \texttt{dataset/images/}).  The additional
212 answer files reference image groups 3 and 4 (participants who selected the third
or fourth candidate image per question), but no image files for groups 3 or 4 have
been provided in the dataset.  Until those images are added, the 212 entries cannot
be scored by the pipeline.  All results in this report are based on the complete
327-sample set for image groups 1--2.
\end{tcolorbox}

% ═══════════════════════════════════════════════════════════════════════════════
\section{Pipeline Architecture}
% ═══════════════════════════════════════════════════════════════════════════════

\subsection{Overview}

The automated scoring pipeline, implemented in \texttt{src/pipeline\_qwen.py}, proceeds
in the following stages for each (image, question, subject-answer) triple:

\begin{enumerate}
  \item \textbf{Translation (He$\to$En):} The Hebrew question is machine-translated to
        English using \textit{Helsinki-NLP/opus-mt-tc-big-he-en}.
  \item \textbf{Concept extraction:} A regex extracts the quoted concept phrase from the
        translated question (e.g., \textit{``The ecosystem is functioning well''}).
  \item \textbf{VLM inference (cached):} \Qwen{} receives the image and the English
        question and produces a systems-thinking expert answer in English.  Answers are
        cached by (image\_path, question\_he), so each of the 20 unique image+question
        pairs requires exactly one VLM call regardless of how many participants answered
        it (see \S\ref{sec:cache}).
  \item \textbf{Subject translation (He$\to$En):} The participant's Hebrew free-text
        justification is translated to English.
  \item \textbf{Similarity scoring:} Cosine similarity is computed between the \vlm{}
        answer and the participant's translated answer using three embedding spaces
        (Table~\ref{tab:embeddings}).
  \item \textbf{Concept baseline:} The cosine similarity between the participant's
        English answer and the bare concept phrase is computed (\csim{}).  Three
        concept-adjusted metrics are then computed: scalar subtraction (\SACs{}),
        vector projection with concept removed from subject only (\SACv{}), and
        bi-projected vector similarity with concept removed from both subject and VLM
        embeddings (\SACvb{}).
  \item \textbf{BERTScore:} Token-level F1 (\bsF{}), Precision, and Recall
        are computed between the \vlm{} reference and the participant's English answer
        using \textit{roberta-large} embeddings.
\end{enumerate}

\begin{table}[H]
\centering
\caption{Embedding models used for similarity scoring.}
\label{tab:embeddings}
\small
\begin{tabular}{lll}
\toprule
\textbf{Metric} & \textbf{Model} & \textbf{Language pair} \\
\midrule
\simHH{}  & paraphrase-multilingual-mpnet-base-v2 & Hebrew VLM answer vs.\ Hebrew subject \\
\simEN{}  & paraphrase-multilingual-mpnet-base-v2 & English VLM answer vs.\ English subject \\
\simLaBSE{} & LaBSE & Hebrew VLM answer vs.\ Hebrew subject \\
\bottomrule
\end{tabular}
\end{table}

\subsection{Design Decisions}

\paragraph{Pre-translating the question.}
Early experiments showed that presenting the Hebrew question directly to \Qwen{} caused
the model to begin its answer with a translation preamble rather than directly engaging
with the image.  Translating the question to English before the VLM call eliminates
this artefact.

\paragraph{VLM answer caching.}
\label{sec:cache}
Because \Qwen{} is a deterministic model (greedy decoding, \texttt{do\_sample=False}),
the same image+question pair always produces the same output.  By caching on
(image\_path, question\_he), the 327-sample run required only \textbf{20 unique VLM
inference calls} (the remaining 307 were served from cache), achieving an effective
$\sim\!15\times$ speedup.

\paragraph{Fallback and degeneracy detection.}
A post-processing function detects degenerate outputs (empty, very short, repetitive
loops, or encoding errors).  In such cases a simpler descriptive question is posed as
a fallback.  In the current dataset, \textbf{zero fallbacks were triggered}, indicating
robust model outputs throughout.

\paragraph{Concept-baseline correction.}
A known confound in free-text ST assessment is \emph{concept echoing}: a participant
who simply restates the concept label (e.g., ``this image shows change being made'')
will score highly on any similarity metric even without genuine visual reasoning.  The
scalar metric \SACs{} corrects for this by subtracting \csim{} from \simEN{}.
The vector metrics \SACv{} and \SACvb{} further remove the concept direction in
embedding space before scoring alignment, providing stricter de-concepted views.
See \S\ref{sec:metrics} for the full formal definition.

% ─────────────────────────────────────────────────────────────────────────────
\subsection{Formal Definition of Scoring Metrics}
\label{sec:metrics}
% ─────────────────────────────────────────────────────────────────────────────

All metrics are defined over a \emph{sample triple} $(i, q, s)$, where $i$ is the
presented image, $q$ is the Hebrew systems-thinking concept question, and $s$ is the
participant's Hebrew free-text justification.
Let $T(\cdot)$ denote the Helsinki He$\to$En machine translation function, and
$\mathbf{v}(i,q)$ the English expert answer produced by the VLM for pair $(i,q)$.

\paragraph{Cosine similarity.}
All embedding-based metrics reduce to cosine similarity over unit-normalised vectors.
For $\mathbf{a}, \mathbf{b} \in \mathbb{R}^{d}$ with $\|\mathbf{a}\|=\|\mathbf{b}\|=1$:
\begin{equation}
  \cos(\mathbf{a}, \mathbf{b})
    = \frac{\mathbf{a} \cdot \mathbf{b}}{\|\mathbf{a}\|\,\|\mathbf{b}\|}
    = \mathbf{a} \cdot \mathbf{b}
  \label{eq:cosine}
\end{equation}
Values lie in $[-1,+1]$; for sentence embeddings the practical range is roughly
$[0, 1]$ for semantically related content.

\paragraph{Primary similarity metrics.}
Let $\phi_{\text{mp}}(\cdot) \in \mathbb{R}^{768}$ be the sentence embedding from
\textit{paraphrase-multilingual-mpnet-base-v2} (normalised), and
$\phi_{\text{LaBSE}}(\cdot) \in \mathbb{R}^{768}$ the embedding from LaBSE.
Let $T^{-1}(\cdot)$ denote the En$\to$He translation model.
The three primary similarity metrics are:

\begin{align}
  \simHH{}(i,q,s)
    &= \cos\!\Bigl(\phi_{\text{mp}}\!\bigl(T^{-1}(\mathbf{v}(i,q))\bigr),\;
                   \phi_{\text{mp}}(s)\Bigr)
    \label{eq:simhh} \\[6pt]
  \simEN{}(i,q,s)
    &= \cos\!\Bigl(\phi_{\text{mp}}(\mathbf{v}(i,q)),\;
                   \phi_{\text{mp}}(T(s))\Bigr)
    \label{eq:simen} \\[6pt]
  \simLaBSE{}(i,q,s)
    &= \cos\!\Bigl(\phi_{\text{LaBSE}}\!\bigl(T^{-1}(\mathbf{v}(i,q))\bigr),\;
                   \phi_{\text{LaBSE}}(s)\Bigr)
    \label{eq:simlabse}
\end{align}

\simHH{} and \simLaBSE{} operate in Hebrew space (both texts in Hebrew);
\simEN{} operates in English space (both texts are machine-translated to English
first), which makes it the most translation-dependent but also the most
semantically explicit metric.

\paragraph{Concept-baseline similarity (\csim{}).}
Let $c(q)$ be the bare English concept phrase extracted from question $q$ by regex
(e.g., \textit{``The ecosystem is functioning well''}).  Then:
\begin{equation}
  \csim{}(q,s)
    = \cos\!\Bigl(\phi_{\text{mp}}(c(q)),\;
                  \phi_{\text{mp}}(T(s))\Bigr)
  \label{eq:csim}
\end{equation}
\csim{} is \emph{image-independent}: it measures how similar a participant's answer is
to the concept label alone, without any reference to the specific image selected.
A high \csim{} means the answer echoes the concept; a low \csim{} (even if \simEN{}
is high) would indicate the participant used image-specific vocabulary not present in
the concept phrase.

\paragraph{Concept-adjusted metrics (project-specific).}
\textbf{1) Scalar baseline adjustment:}
\begin{equation}
      \boxed{\SACs{}(i,q,s)
            \;=\; \simEN{}(i,q,s)
            \;-\; \csim{}(q,s)}
      \label{eq:sac_scalar}
\end{equation}

\textbf{2) Subject-only projected vector metric:}
\begin{equation}
      \boxed{\SACv{}(i,q,s)
            \;=\;
            \cos\!\Bigl(\widetilde{\mathbf{u}}_{\perp c},\; \widetilde{\mathbf{v}}\Bigr)}
      \label{eq:sac_vec}
\end{equation}
where
\begin{equation}
      \widetilde{\mathbf{u}}_{\perp c}
      =
      \frac{\widetilde{\mathbf{u}} - (\widetilde{\mathbf{u}}\!\cdot\!\widetilde{\mathbf{c}})\widetilde{\mathbf{c}}}
                   {\left\|\widetilde{\mathbf{u}} - (\widetilde{\mathbf{u}}\!\cdot\!\widetilde{\mathbf{c}})\widetilde{\mathbf{c}}\right\|}
\end{equation}
and $\widetilde{\mathbf{u}}, \widetilde{\mathbf{v}}, \widetilde{\mathbf{c}}$ are the
unit-normalized embeddings of subject answer, VLM answer, and concept phrase.

\textbf{3) Bi-projected vector metric (new):}
\begin{equation}
      \boxed{\SACvb{}(i,q,s)
            \;=\;
            \cos\!\Bigl(\widetilde{\mathbf{u}}_{\perp c},\; \widetilde{\mathbf{v}}_{\perp c}\Bigr)}
      \label{eq:sac_vec_both}
\end{equation}
with
\begin{equation}
      \widetilde{\mathbf{v}}_{\perp c}
      =
      \frac{\widetilde{\mathbf{v}} - (\widetilde{\mathbf{v}}\!\cdot\!\widetilde{\mathbf{c}})\widetilde{\mathbf{c}}}
                   {\left\|\widetilde{\mathbf{v}} - (\widetilde{\mathbf{v}}\!\cdot\!\widetilde{\mathbf{c}})\widetilde{\mathbf{c}}\right\|}
\end{equation}

Interpretation:
\begin{itemize}[noitemsep]
      \item $\SACs{} > 0$ --- the answer is closer to the VLM's image-specific interpretation
                        than to the bare concept, implying \textbf{visual grounding}.
      \item $\SACs{} \approx 0$ --- the answer adds no image-specific information; the
                        participant is writing about the concept, not the image.
      \item $\SACs{} < 0$ --- the participant's language is semantically more similar to the
                        concept label than to the VLM reference; a strong indicator of
                        \textbf{concept echoing} with minimal visual engagement.
      \item $\SACv{}$ and $\SACvb{}$ are cosine similarities in de-concepted subspaces,
                        so both lie in $[-1,+1]$; in this dataset they are empirically positive for all
                        samples.
\end{itemize}

Grounding rates are defined as fractions of samples above zero per metric:
\begin{equation}
      \gamma_s
            = \frac{1}{N}
                  \sum_{n=1}^{N}
                  \mathbf{1}\!\bigl[\SACs{}(i_n, q_n, s_n) > 0\bigr],
      \quad
      \gamma_v
            = \frac{1}{N}\sum_{n=1}^{N}\mathbf{1}\!\bigl[\SACv{}(i_n, q_n, s_n) > 0\bigr],
      \quad
      \gamma_{vb}
            = \frac{1}{N}\sum_{n=1}^{N}\mathbf{1}\!\bigl[\SACvb{}(i_n, q_n, s_n) > 0\bigr]
      \label{eq:gamma_multi}
\end{equation}
In the current 327-sample dataset:
$\gamma_s=0.621$ (62.1\%), $\gamma_v=1.000$ (100.0\%), and $\gamma_{vb}=1.000$ (100.0\%).

Empirical ranges in this dataset are:
$\SACs{}\in[-0.516, 0.537]$,
$\SACv{}\in[0.006, 0.765]$,
$\SACvb{}\in[0.007, 0.799]$
(Table~\ref{tab:global_stats}).

\paragraph{BERTScore.}
BERTScore \citep{bertscore2020} operates at the token level rather than the
sentence level, providing a complementary perspective.
Let $\{r_k\}$ be the tokens of the VLM reference $\mathbf{v}(i,q)$ and $\{h_j\}$
the tokens of the participant's English translation $T(s)$, with contextual embeddings
$\mathbf{e}_{r_k}$, $\mathbf{e}_{h_j}$ from \textit{roberta-large}:
\begin{equation}
  P_{\text{BERT}} = \frac{1}{|H|}\sum_{j}
    \max_{k}\,\cos(\mathbf{e}_{h_j},\, \mathbf{e}_{r_k}),
  \qquad
  R_{\text{BERT}} = \frac{1}{|R|}\sum_{k}
    \max_{j}\,\cos(\mathbf{e}_{r_k},\, \mathbf{e}_{h_j})
  \label{eq:bs_pr}
\end{equation}
\begin{equation}
  \bsF{}
    = \frac{2\,P_{\text{BERT}}\cdot R_{\text{BERT}}}
           {P_{\text{BERT}} + R_{\text{BERT}}}
  \label{eq:bertscore}
\end{equation}
The consistent pattern $R_{\text{BERT}} > P_{\text{BERT}}$ observed here (means 0.867
vs.\ 0.825) is expected: participant answers tend to contain vocabulary subsumed by the
longer VLM reference, but miss some of the VLM's image-specific tokens --- hence high
recall with lower precision.

\subsection{Hardware and Models}

\begin{itemize}[noitemsep]
  \item \textbf{VLM:} \Qwen{} (4-bit NF4 quantisation via \texttt{bitsandbytes}),
        executed on GPU (CUDA).
  \item \textbf{Translation:} Helsinki opus-mt-tc-big-he-en (He$\to$En) and
        opus-mt-en-he (En$\to$He).
  \item \textbf{Embeddings:} paraphrase-multilingual-mpnet-base-v2 and LaBSE via
        \textit{sentence-transformers}.
  \item \textbf{BERTScore:} \textit{roberta-large}, CPU evaluation.
\end{itemize}

% ═══════════════════════════════════════════════════════════════════════════════
\section{Preliminary Results (N = 327)}
% ═══════════════════════════════════════════════════════════════════════════════

\subsection{Similarity Score Distributions}

Table~\ref{tab:global_stats} reports global descriptive statistics for all similarity
metrics across the 327 fully-scored samples.

\begin{table}[H]
\centering
\caption{Descriptive statistics for all similarity and quality metrics ($N=327$).}
\label{tab:global_stats}
\small
\begin{tabular}{lrrrrrr}
\toprule
\textbf{Metric} & \textbf{Mean} & \textbf{SD} & \textbf{Min} & \textbf{Median} & \textbf{Max} \\
\midrule
\simHH{}         & 0.535 & 0.136 & 0.153 & 0.543 & 0.898 \\
\simEN{}         & 0.494 & 0.166 & 0.015 & 0.505 & 0.828 \\
\simLaBSE{}      & 0.373 & 0.129 & 0.016 & 0.369 & 0.809 \\
\csim{}          & 0.419 & 0.170 & 0.017 & 0.412 & 0.826 \\
\SACs{}          & 0.075 & 0.197 & $-$0.516 & 0.059 & 0.537 \\
\SACv{}          & 0.350 & 0.174 & 0.006 & 0.354 & 0.765 \\
\SACvb{}         & 0.389 & 0.191 & 0.007 & 0.394 & 0.799 \\
\bsF{} (Recall)  & 0.867 & 0.021 & 0.809 & 0.868 & 0.943 \\
\bsF{} (Precision) & 0.825 & 0.017 & 0.773 & 0.825 & 0.895 \\
\bsF{} (F1)      & 0.845 & 0.016 & 0.791 & 0.845 & 0.889 \\
Subject word count & 12.4 & 11.1 & 1 & 9.0 & 75 \\
\bottomrule
\end{tabular}
\end{table}

Several observations are noteworthy:

\begin{itemize}
  \item \simHH{} (He--He multilingual embeddings) yields the highest mean
        ($\bar{x}=0.535$), likely because both texts share Hebrew vocabulary
        and the multilingual encoder has reduced cross-lingual alignment noise.
  \item \simEN{} is slightly lower ($\bar{x}=0.494$) due to translation noise
        introduced by the machine-translation step, but has the widest range
        ($[0.015,\,0.828]$), making it the most discriminating metric across samples.
  \item \simLaBSE{} is systematically lower ($\bar{x}=0.373$), consistent with
        LaBSE's design for document-level rather than short-phrase alignment.
  \item The concept-baseline \csim{} has a mean of 0.419 -- close to \simEN{} -- 
        confirming that a substantial fraction of similarity is attributable to
        concept echoing alone.
  \item \SACs{} has a positive median (0.059) and mean (0.075), but also a substantial
        negative tail ($-0.516$), indicating heterogeneous concept-echoing behavior.
  \item \SACv{} and \SACvb{} are both strongly positive on average (0.350 and 0.389,
        respectively), with \SACvb{} consistently slightly higher because concept
        direction is removed from both subject and VLM representations.
\end{itemize}

\paragraph{Paired comparison of \simHH{} vs.\ \simEN{}.}
A paired two-tailed $t$-test confirms that \simHH{} is significantly higher than
\simEN{} ($t(326)=6.43$, $p=4.66\times10^{-10}$), consistent with the hypothesis that
operating in the same language space inflates similarity.  The Pearson correlation
between the two metrics is $r=0.720$ ($p<10^{-50}$), indicating strong but imperfect
agreement.

\subsection{Concept Baseline and Grounding Analysis}

Of the 327 samples, \textbf{62.1\%} had $\SACs{}>0$, while \textbf{100\%} had
$\SACv{}>0$ and \textbf{100\%} had $\SACvb{}>0$.
The scalar metric is therefore the only one that produces negative values in this
dataset and remains the primary diagnostic for explicit concept-echoing failures.

Table~\ref{tab:per_concept} presents per-concept statistics.

\begin{table}[H]
\centering
\caption{Per-concept similarity and grounding statistics ($N=327$, partial data).}
\label{tab:per_concept}
\small
\begin{tabular}{clrrrr}
\toprule
\textbf{Q} & \textbf{Concept (abbreviated)} & \textbf{N} & \textbf{\simEN{} (M±SD)} & \textbf{\SACs{} mean} & \textbf{\% grounded (\SACs{}$>0$)} \\
\midrule
Q1  & Work completed successfully              & 14 & $0.425 \pm 0.136$ & $+0.072$ & 71\% \\
Q2  & The decision has been made               & 57 & $0.516 \pm 0.156$ & $+0.089$ & 63\% \\
Q3  & Child learned an important lesson        & 33 & $0.471 \pm 0.149$ & $+0.004$ & 45\% \\
Q4  & Ecosystem is functioning well            & 39 & $0.490 \pm 0.214$ & $+0.143$ & 67\% \\
Q5  & Similar systems                          & 27 & $0.523 \pm 0.195$ & $+0.159$ & 78\% \\
Q6  & Group advanced towards goal              & 31 & $0.558 \pm 0.178$ & $+0.138$ & 77\% \\
Q7  & Understanding reality affects what I see & 38 & $0.470 \pm 0.141$ & $-0.008$ & 50\% \\
Q8  & Seeing things differently                & 13 & $0.534 \pm 0.088$ & $+0.014$ & 46\% \\
Q9  & Man affects the world around him         & 36 & $0.469 \pm 0.165$ & $+0.091$ & 64\% \\
Q10 & Change is made                           & 39 & $0.471 \pm 0.149$ & $+0.029$ & 59\% \\
\midrule
\textbf{Total} & & \textbf{327} & $\mathbf{0.494 \pm 0.166}$ & $\mathbf{+0.075}$ & \textbf{62\%} \\
\bottomrule
\end{tabular}
\end{table}

Notable patterns:

\begin{itemize}
  \item \textbf{Q6} (``Group advanced towards goal'') achieves the highest mean \simEN{}
        ($0.558$) and second-highest grounding rate (77\%), suggesting participants
        gave descriptively rich, image-specific justifications.
  \item \textbf{Q5} (``Similar systems'') shows the highest grounding rate (78\%) and
        second-highest \SACs{} ($+0.159$), despite the concept being abstract -- a
        possible sign that the two associated images are unambiguous.
  \item \textbf{Q3} (``Child learned an important lesson'') has the lowest grounding
        rate (45\%) and near-zero \SACs{} ($+0.004$), consistent with participants
        simply asserting the concept label rather than reasoning visually.
  \item \textbf{Q7} (``Understanding reality affects what I see'') is the only concept
        with a \emph{negative} mean \SACs{} ($-0.008$) and only 50\% grounding -- the
        high abstraction of Q7 may cause participants to rely heavily on meta-reflective
        vocabulary that echoes the concept more than it reflects the image.
\end{itemize}

\subsection{Answer Length and Similarity}

The length of a subject's Hebrew answer (word count) is positively correlated with
\simEN{} (Spearman $\rho = 0.40$, $p < 10^{-13}$).  Table~\ref{tab:length_brackets}
shows that this relationship holds monotonically across three length brackets.

\begin{table}[H]
\centering
\caption{Mean similarity by subject answer length.}
\label{tab:length_brackets}
\small
\begin{tabular}{lrr}
\toprule
\textbf{Length bracket} & \textbf{Mean \simEN{}} & \textbf{Mean \SACs{}} \\
\midrule
Short ($1$--$5$ words)  & 0.387 & 0.044 \\
Medium ($6$--$20$ words)& 0.512 & 0.076 \\
Long ($>20$ words)      & 0.583 & 0.121 \\
\bottomrule
\end{tabular}
\end{table}

The increasing \SACs{} across length brackets (0.044 $\to$ 0.076 $\to$ 0.121) is
particularly informative: longer responses not only score higher overall, but
they do so by going \emph{beyond} the concept, rather than merely being verbose
repetitions of the concept phrase.

\subsection{BERTScore Analysis}

BERTScore F1 (\bsF{}) has a narrower range ($[0.791, 0.889]$, mean $0.845$, SD $0.016$)
compared to cosine similarity, indicating lower discriminative power for this task.
The Pearson correlation between \bsF{} and \simEN{} is $r=0.584$ ($p<10^{-30}$),
confirming moderate agreement but suggesting the two metrics capture partially
independent aspects of response quality.

BERTScore recall (mean $0.867$) consistently exceeds precision (mean $0.825$), which
makes sense for this task: participant answers tend to contain substrings of the VLM
reference vocabulary, but the VLM answers also contain information not present in the
shorter, more telegraphic participant answers.

\subsection{Metric Inter-Correlations}
\label{sec:correlations}

The strongest Spearman correlations among all numeric metrics are:

\begin{enumerate}[noitemsep]
  \item $\simHH{}$ vs.\ \simEN{}: $\rho = 0.72$ -- same metric, different language
        spaces.
  \item \simEN{} vs.\ \bsF{}: $\rho = 0.58$ -- moderate agreement between embedding
        cosine and token-level overlap.
  \item \SACv{} vs.\ \SACvb{}: $\rho \approx 0.99$ -- both vector metrics are nearly
        collinear, with \SACvb{} systematically slightly higher.
  \item \simEN{} vs.\ \SACvb{}: $\rho \approx 0.94$ -- de-concepted and raw alignment
        remain tightly coupled in this dataset.
  \item Subject word count vs.\ \simEN{}: $\rho = 0.40$ -- length predicts quality
        but is not sufficient.
\end{enumerate}

% ═══════════════════════════════════════════════════════════════════════════════
\section{Discussion}
% ═══════════════════════════════════════════════════════════════════════════════

\subsection{Validity of the \vlm{}-as-Reference Approach}

The pipeline produces a graded, continuous score for open-ended Hebrew reasoning about
images, without requiring human annotation.  Several lines of evidence support its
construct validity:

\begin{enumerate}
  \item The score is sensitive to response length in the expected direction, and the
        increase in \SACs{} with length (Table~\ref{tab:length_brackets}) indicates it
        responds to content richness, not mere length inflation.
  \item The concept-baseline correction (\SACs{}) successfully identifies the distinction
        between concept-echoing responses and image-grounded responses.  Q3 and Q7 --
        both abstract, reflexive concepts -- cluster at low \SACs{}, while Q5 and Q6 --
        perceptually concrete -- cluster at high \SACs{}, which aligns with theoretical
        expectations.
  \item Zero fallbacks were triggered, suggesting \Qwen{} produces consistent,
        coherent expert answers across the 20 unique image+question pairs.
\end{enumerate}

\subsection{Limitations}

\begin{itemize}
  \item \textbf{Translation noise.} Both the question pre-translation and the subject
        answer translation introduce errors.  Hebrew is morphologically complex, and
        machine translation at the short-phrase level (1--20 words) is prone to 
        mis-translating subtle reasoning markers.  A future direction is to use a
        Hebrew-native embedding model directly (e.g., a Hebrew BERT) to bypass this.
  \item \textbf{Single VLM reference.} All samples for the same (image, question) pair
        receive an identical VLM reference.  Inter-image variability in the concept
        is collapsed: participants who chose image 1 for Q6 are compared against the
        same reference regardless of their individual interpretation.  A per-participant
        VLM call with personalised prompting could address this.
  \item \textbf{Partial dataset.} Results cover 327 of 539 answer files (image
        groups 1--2 only).  The 212 remaining answer files reference image groups 3--4,
        for which no image files exist in the dataset at the time of writing.
        If those images are later added, the pipeline can be resumed and the analysis
        extended to all 539 samples.
  \item \textbf{Missing-text entries.} The missing-pairs audit (167 entries with
        \texttt{missing\_text}, \textasciitilde{}30 with \texttt{missing\_option})
        suggests a non-trivial dropout in the questionnaire, concentrated in certain
        subjects (e.g., subjects 1, 2, 3, 6, 51, 65, 66).  These entries are excluded
        from the analysis but may represent a systematic non-response bias.
\end{itemize}

\subsection{Implications for ST Assessment}

The \SACs{} metric operationalises a key theoretical distinction in ST measurement:
the difference between \emph{recognising} a concept (selecting the right image) and
\emph{applying} it with visual grounding (writing a substantive justification).
The finding that 38\% of participants score $\SACs{} \le 0$ while still having selected
a plausible image suggests that image selection alone overestimates ST competence.
This reinforces the importance of free-text justification in ST assessment batteries.

% ═══════════════════════════════════════════════════════════════════════════════
\section{Methods Summary}
% ═══════════════════════════════════════════════════════════════════════════════

\subsection{Software and Reproducibility}

The full pipeline is implemented in Python 3.10 and is available in the project
repository.  Key dependencies:

\begin{lstlisting}
transformers>=4.40      # Qwen2-VL, Marian translation
sentence-transformers   # mpnet, LaBSE embeddings
bert-score              # BERTScorer (roberta-large)
bitsandbytes            # 4-bit NF4 quantisation
qwen-vl-utils           # vision input preprocessing
torch>=2.1              # CUDA inference
\end{lstlisting}

The pipeline can be resumed from a partial results file via:

\begin{lstlisting}[language=bash]
python -m src.pipeline_qwen \
    --dataset dataset \
    --resume results/vlm_results_qwen_20260227_152809.json
\end{lstlisting}

\subsection{Scoring Metrics Reference}

Full formal definitions with equations are in \S\ref{sec:metrics}.  Brief reference:

\begin{description}[leftmargin=3.2cm, labelwidth=3cm]
  \item[\simHH{}]       Eq.~\eqref{eq:simhh}: cosine, multilingual mpnet,
                        He VLM answer vs.\ He subject answer.
  \item[\simEN{}]       Eq.~\eqref{eq:simen}: cosine, multilingual mpnet,
                        En VLM answer vs.\ En subject answer (both machine-translated).
  \item[\simLaBSE{}]    Eq.~\eqref{eq:simlabse}: cosine, LaBSE,
                        He VLM answer vs.\ He subject answer.
  \item[\csim{}]        Eq.~\eqref{eq:csim}: cosine between En subject answer
                        and bare English concept phrase (image-independent baseline).
      \item[\SACs{}]        Eq.~\eqref{eq:sac_scalar}: scalar baseline subtraction,
                                                                        $\simEN{} - \csim{}$.
      \item[\SACv{}]        Eq.~\eqref{eq:sac_vec}: concept removed from subject
                                                                        embedding, then cosine to VLM embedding.
      \item[\SACvb{}]       Eq.~\eqref{eq:sac_vec_both}: concept removed from both
                                                                        subject and VLM embeddings before cosine.
  \item[\bsF{}]         Eqs.~\eqref{eq:bs_pr}--\eqref{eq:bertscore}: BERTScore F1
                        (roberta-large), En VLM answer vs.\ En subject.
\end{description}

% ═══════════════════════════════════════════════════════════════════════════════
\section{Next Steps}
% ═══════════════════════════════════════════════════════════════════════════════

\begin{enumerate}
  \item \textbf{Provide images for groups 3--4.} Adding the 20 missing image files
        (\texttt{<q>\_3.png}, \texttt{<q>\_4.png} for $q=1\ldots10$) to
        \texttt{dataset/images/} would make the remaining 212 answer files scoreable.
        The pipeline supports resuming from a checkpoint
        (\texttt{--resume} flag), so only new samples would be processed.
  \item \textbf{Human validation.} Annotate a stratified sample ($n \approx 60$) of
        responses across the score distribution with expert human ratings to establish
      criterion validity of \SACs{}, \SACv{}, \SACvb{}, and \simEN{}.
  \item \textbf{Per-subject profiling.} Aggregate scores to the subject level to
        produce individual ST profiles, and correlate with demographic and
        background variables.
  \item \textbf{Predictive modelling.} Explore whether word count, concept, and
        image variables jointly predict scores in a mixed-effects model.
  \item \textbf{Hebrew-native embeddings.} Test Hebrew BERT / alephBERT as an
        alternative embedding backbone to reduce translation noise.
  \item \textbf{Missing-data strategy.} Investigate whether participants with
        \texttt{missing\_text} entries differ systematically from those with complete
        data (non-ignorable vs.\ ignorable missingness).
\end{enumerate}

% ═══════════════════════════════════════════════════════════════════════════════
\section*{Acknowledgements}
% ═══════════════════════════════════════════════════════════════════════════════

Pipeline developed using open-weight models from Qwen (Alibaba), Helsinki NLP,
and sentence-transformers.  Compute resources: local CUDA GPU workstation.

% ───────────────────────────────────────────────────────────────────────────────
\begin{thebibliography}{9}

\bibitem[Senge(1990)]{senge1990fifth}
Senge, P.~M. (1990).
\newblock \textit{The Fifth Discipline: The Art and Practice of the Learning Organization}.
\newblock Doubleday/Currency.

\bibitem[Assaraf \& Orion(2005)]{assaraf2005}
Assaraf, O.~B.-Z., \& Orion, N. (2005).
\newblock Development of system thinking skills in the context of earth system education.
\newblock \textit{Journal of Research in Science Teaching}, 42(5), 518--560.

\bibitem[Hmelo-Silver et al.(2007)]{hmelo2007}
Hmelo-Silver, C.~E., Marathe, S., \& Liu, L. (2007).
\newblock Fish swim, rocks sit, and lungs breathe: Expert-novice understanding of
  complex systems.
\newblock \textit{Journal of the Learning Sciences}, 16(3), 307--331.

\bibitem[Zhang et al.(2020)]{bertscore2020}
Zhang, T., Kishore, V., Wu, F., Weinberger, K.~Q., \& Artzi, Y. (2020).
\newblock BERTScore: Evaluating text generation with BERT.
\newblock In \textit{ICLR 2020}.

\bibitem[Wang et al.(2024)]{qwen2vl2024}
Wang, P., et al. (2024).
\newblock Qwen2-VL: Enhancing vision-language model's perception of the world at any resolution.
\newblock \textit{arXiv:2409.12191}.

\bibitem[Feng et al.(2022)]{labse2022}
Feng, F., et al. (2022).
\newblock Language-agnostic BERT sentence embedding.
\newblock In \textit{Proceedings of ACL 2022}.

\end{thebibliography}

\end{document}
